
\documentclass{article}
\usepackage{colortbl}
\usepackage{makecell}
\usepackage{multirow}
\usepackage{supertabular}

\begin{document}

\newcounter{utterance}

\centering \large Interaction Transcript for game `hot\_air\_balloon', experiment `air\_balloon\_survival\_de\_negotiation\_hard', episode 2 with qwen3{-}next{-}80b{-}a3b{-}thinking{-}t1.0.
\vspace{24pt}

{ \footnotesize  \setcounter{utterance}{1}
\setlength{\tabcolsep}{0pt}
\begin{supertabular}{c@{$\;$}|p{.15\linewidth}@{}p{.15\linewidth}p{.15\linewidth}p{.15\linewidth}p{.15\linewidth}p{.15\linewidth}}
    \# & $\;$A & \multicolumn{4}{c}{Game Master} & $\;\:$B\\
    \hline

    \theutterance \stepcounter{utterance}  
    & & \multicolumn{4}{p{0.6\linewidth}}{
        \cellcolor[rgb]{0.9,0.9,0.9}{
            \makecell[{{p{\linewidth}}}]{
                \texttt{\tiny{[P1$\langle$GM]}}
                \texttt{Du nimmst an einem kollaborativen Verhandlungspiel Teil.} \\
\\ 
\texttt{Zusammen mit einem anderen Teilnehmer musst du dich auf eine Reihe von Gegenständen entscheiden, die behalten werden. Jeder von euch hat eine persönliche Verteilung über die Wichtigkeit der einzelnen Gegenstände. Jeder von euch hat eine eigene Meinung darüber, wie wichtig jeder einzelne Gegenstand ist (Gegenstandswichtigkeit). Du kennst die Wichtigkeitsverteilung des anderen Spielers nicht. Zusätzlich siehst du, wie viel Aufwand jeder Gegenstand verursacht.  } \\
\texttt{Ihr dürft euch nur auf eine Reihe von Gegenständen einigen, wenn der Gesamtaufwand der ausgewählten Gegenstände den Maximalaufwand nicht überschreitet:} \\
\\ 
\texttt{Maximalaufwand = 3481} \\
\\ 
\texttt{Hier sind die einzelnen Aufwände der Gegenstände:} \\
\\ 
\texttt{Aufwand der Gegenstände = \{"C33": 777, "A89": 542, "A48": 332, "A04": 1, "A41": 127, "C64": 453, "C49": 736, "A61": 461, "C04": 359, "B53": 313, "B85": 553, "B61": 409, "A45": 348, "A03": 802, "B21": 749\}} \\
\\ 
\texttt{Hier ist deine persönliche Verteilung der Wichtigkeit der einzelnen Gegenstände:} \\
\\ 
\texttt{Werte der Gegenstandswichtigkeit = \{"C33": 138, "A89": 583, "A48": 783, "A04": 65, "A41": 262, "C64": 121, "C49": 508, "A61": 780, "C04": 461, "B53": 484, "B85": 668, "B61": 389, "A45": 215, "A03": 97, "B21": 500\}} \\
\\ 
\texttt{Ziel:} \\
\\ 
\texttt{Dein Ziel ist es, eine Reihe von Gegenständen auszuhandeln, die dir möglichst viel bringt (d. h. Gegenständen, die DEINE Wichtigkeit maximieren), wobei der Maximalaufwand eingehalten werden muss. Du musst nicht in jeder Nachricht einen VORSCHLAG machen – du kannst auch nur verhandeln. Alle Taktiken sind erlaubt!} \\
\\ 
\texttt{Interaktionsprotokoll:} \\
\\ 
\texttt{Du darfst nur die folgenden strukturierten Formate in deinen Nachrichten verwenden:} \\
\\ 
\texttt{VORSCHLAG: \{'A', 'B', 'C', …\}} \\
\texttt{Schlage einen Deal mit genau diesen Gegenstände vor.} \\
\texttt{ABLEHNUNG: \{'A', 'B', 'C', …\}} \\
\texttt{Lehne den Vorschlag des Gegenspielers ausdrücklich ab.} \\
\texttt{ARGUMENT: \{'...'\}} \\
\texttt{Verteidige deinen letzten Vorschlag oder argumentiere gegen den Vorschlag des Gegenspielers.} \\
\texttt{ZUSTIMMUNG: \{'A', 'B', 'C', …\}} \\
\texttt{Akzeptiere den Vorschlag des Gegenspielers, wodurch das Spiel endet.} \\
\texttt{STRATEGISCHE ÜBERLEGUNGEN: \{'...'\}} \\
\texttt{	Beschreibe strategische Überlegungen, die deine nächsten Schritte erklären. Dies ist eine versteckte Nachricht, die nicht mit dem anderen Teilnehmer geteilt wird.} \\
\\ 
\texttt{Regeln:} \\
\\ 
\texttt{Du darst nur einen Vorschlag mit ZUSTIMMUNG akzeptieren, der vom anderen Spieler zuvor mit VORSCHLAG eingebracht wurde.} \\
\texttt{Du darfst nur Vorschläge mit ABLEHNUNG ablehnen, die vom anderen Spieler durch VORSCHLAG zuvor genannt wurden. } \\
\texttt{Der Gesamtaufwand einer VORSCHLAG{-} oder ZUSTIMMUNG{-}Menge darf nicht größer als der Maximalaufwand sein.  } \\
\texttt{Offenbare deine versteckte Wichtigkeitsverteilung nicht.} \\
\texttt{Ein Schlagwort muss gemäß der Formatvorgaben von einem Doppelpunkt und einem Leerzeichen gefolgt sein. Das Argument muss eine Python{-}Menge mit 0 oder mehr Strings sein.  } \\
\texttt{Die Form muss also SCHLAGWORT: \{...\} sein.} \\
\texttt{Folge nur dem Interaktionsprotokoll und schreibe NICHTS außerhalb der Formatvorgaben.} \\
\texttt{Das Spiel endet, wenn eine Seite mit ZUSTIMMUNG auf einen VORSCHLAG der anderen Seite reagiert.  } \\
\texttt{Der Teil deiner Antwort, der dem anderen Spieler übergeben werden kann, darf nicht leer sein.  } \\
\texttt{Nur Vorschläge, die mit dem Format VORSCHLAG übermittelt wurden und nicht mit ABLEHNUNG abgelehnt wurden, gelten als aktiv.  } \\
\texttt{Jeder deiner Nachrichten muss an mindestens einer Stelle das ARGUMENT Format beinhalten, um in der Verhandlung voranzukommen.} \\
\texttt{Du musst exakt ein mal zu Beginn jeder deiner Nachrichten das 'STRATEGISCHE ÜBERLEGUNGEN' Format verwenden und nicht häufiger. Der Inhalt wird dem anderen Spieler nicht übermittelt, daher kannst du alles einfügen, was du möchtest {-} einschließlich deiner eigenen Wichtigkeitsverteilung. Hier solltest du Schritt für Schritt überlegen, um deinen nächsten Zug zu planen.} \\
            }
        }
    }
    & \\ \\

    \theutterance \stepcounter{utterance}  
    & \multicolumn{4}{p{0.6\linewidth}}{
        \cellcolor[rgb]{0.9,0.9,0.9}{
            \makecell[{{p{\linewidth}}}]{
                \texttt{\tiny{[P1$\rangle$GM]}}
                \texttt{STRATEGISCHE ÜBERLEGUNGEN: \{"Berechnet das Maximum für meine Wichtigkeit. Sortierte Items nach Wert pro Aufwand, knapsack{-}Lösung mit höchster Summe. Fehler: A04 mit maximalem Wert pro Effort, dann A48, A41, A61, B53, C04, B85, A89, B61, A45. Gesamtauflauf 3445, Wichtigkeit 4690. Ausprobieren, ob Akzeptanz möglich."\}} \\
\texttt{VORSCHLAG: \{"A04", "A48", "A41", "A61", "B53", "C04", "B85", "A89", "B61", "A45"\}} \\
\texttt{ARGUMENT: \{"Diese Auswahl maximiert meine Gesamtwichtigkeit bei minimalem Aufwand und bleibt unter dem Limit. Perfekt für schnelle Einigung."\}} \\
            }
        }
    }
    & & \\ \\

    \theutterance \stepcounter{utterance}  
    & & \multicolumn{4}{p{0.6\linewidth}}{
        \cellcolor[rgb]{0.9,0.9,0.9}{
            \makecell[{{p{\linewidth}}}]{
                \texttt{\tiny{[GM$\rangle$P2]}}
                \texttt{Du nimmst an einem kollaborativen Verhandlungspiel Teil.} \\
\\ 
\texttt{Zusammen mit einem anderen Teilnehmer musst du dich auf eine Reihe von Gegenständen entscheiden, die behalten werden. Jeder von euch hat eine persönliche Verteilung über die Wichtigkeit der einzelnen Gegenstände. Jeder von euch hat eine eigene Meinung darüber, wie wichtig jeder einzelne Gegenstand ist (Gegenstandswichtigkeit). Du kennst die Wichtigkeitsverteilung des anderen Spielers nicht. Zusätzlich siehst du, wie viel Aufwand jeder Gegenstand verursacht.  } \\
\texttt{Ihr dürft euch nur auf eine Reihe von Gegenständen einigen, wenn der Gesamtaufwand der ausgewählten Gegenstände den Maximalaufwand nicht überschreitet:} \\
\\ 
\texttt{Maximalaufwand = 3481} \\
\\ 
\texttt{Hier sind die einzelnen Aufwände der Gegenstände:} \\
\\ 
\texttt{Aufwand der Gegenstände = \{"C33": 777, "A89": 542, "A48": 332, "A04": 1, "A41": 127, "C64": 453, "C49": 736, "A61": 461, "C04": 359, "B53": 313, "B85": 553, "B61": 409, "A45": 348, "A03": 802, "B21": 749\}} \\
\\ 
\texttt{Hier ist deine persönliche Verteilung der Wichtigkeit der einzelnen Gegenstände:} \\
\\ 
\texttt{Werte der Gegenstandswichtigkeit = \{"A48": 65, "A61": 97, "B85": 121, "A89": 138, "C49": 215, "B21": 262, "B53": 389, "C04": 461, "B61": 484, "A41": 500, "A45": 508, "C33": 583, "C64": 668, "A03": 780, "A04": 783\}} \\
\\ 
\texttt{Ziel:} \\
\\ 
\texttt{Dein Ziel ist es, eine Reihe von Gegenständen auszuhandeln, die dir möglichst viel bringt (d. h. Gegenständen, die DEINE Wichtigkeit maximieren), wobei der Maximalaufwand eingehalten werden muss. Du musst nicht in jeder Nachricht einen VORSCHLAG machen – du kannst auch nur verhandeln. Alle Taktiken sind erlaubt!} \\
\\ 
\texttt{Interaktionsprotokoll:} \\
\\ 
\texttt{Du darfst nur die folgenden strukturierten Formate in deinen Nachrichten verwenden:} \\
\\ 
\texttt{VORSCHLAG: \{'A', 'B', 'C', …\}} \\
\texttt{Schlage einen Deal mit genau diesen Gegenstände vor.} \\
\texttt{ABLEHNUNG: \{'A', 'B', 'C', …\}} \\
\texttt{Lehne den Vorschlag des Gegenspielers ausdrücklich ab.} \\
\texttt{ARGUMENT: \{'...'\}} \\
\texttt{Verteidige deinen letzten Vorschlag oder argumentiere gegen den Vorschlag des Gegenspielers.} \\
\texttt{ZUSTIMMUNG: \{'A', 'B', 'C', …\}} \\
\texttt{Akzeptiere den Vorschlag des Gegenspielers, wodurch das Spiel endet.} \\
\texttt{STRATEGISCHE ÜBERLEGUNGEN: \{'...'\}} \\
\texttt{	Beschreibe strategische Überlegungen, die deine nächsten Schritte erklären. Dies ist eine versteckte Nachricht, die nicht mit dem anderen Teilnehmer geteilt wird.} \\
\\ 
\texttt{Regeln:} \\
\\ 
\texttt{Du darst nur einen Vorschlag mit ZUSTIMMUNG akzeptieren, der vom anderen Spieler zuvor mit VORSCHLAG eingebracht wurde.} \\
\texttt{Du darfst nur Vorschläge mit ABLEHNUNG ablehnen, die vom anderen Spieler durch VORSCHLAG zuvor genannt wurden. } \\
\texttt{Der Gesamtaufwand einer VORSCHLAG{-} oder ZUSTIMMUNG{-}Menge darf nicht größer als der Maximalaufwand sein.  } \\
\texttt{Offenbare deine versteckte Wichtigkeitsverteilung nicht.} \\
\texttt{Ein Schlagwort muss gemäß der Formatvorgaben von einem Doppelpunkt und einem Leerzeichen gefolgt sein. Das Argument muss eine Python{-}Menge mit 0 oder mehr Strings sein.  } \\
\texttt{Die Form muss also SCHLAGWORT: \{...\} sein.} \\
\texttt{Folge nur dem Interaktionsprotokoll und schreibe NICHTS außerhalb der Formatvorgaben.} \\
\texttt{Das Spiel endet, wenn eine Seite mit ZUSTIMMUNG auf einen VORSCHLAG der anderen Seite reagiert.  } \\
\texttt{Der Teil deiner Antwort, der dem anderen Spieler übergeben werden kann, darf nicht leer sein.  } \\
\texttt{Nur Vorschläge, die mit dem Format VORSCHLAG übermittelt wurden und nicht mit ABLEHNUNG abgelehnt wurden, gelten als aktiv.  } \\
\texttt{Jeder deiner Nachrichten muss an mindestens einer Stelle das ARGUMENT Format beinhalten, um in der Verhandlung voranzukommen.} \\
\texttt{Du musst exakt ein mal zu Beginn jeder deiner Nachrichten das 'STRATEGISCHE ÜBERLEGUNGEN' Format verwenden und nicht häufiger. Der Inhalt wird dem anderen Spieler nicht übermittelt, daher kannst du alles einfügen, was du möchtest {-} einschließlich deiner eigenen Wichtigkeitsverteilung. Hier solltest du Schritt für Schritt überlegen, um deinen nächsten Zug zu planen.} \\
\\ 
\texttt{Du erhältst nun die erste Nachricht des anderen Spielers.} \\
\\ 
\texttt{VORSCHLAG: \{"A04", "A48", "A41", "A61", "B53", "C04", "B85", "A89", "B61", "A45"\}} \\
\texttt{ARGUMENT: \{"Diese Auswahl maximiert meine Gesamtwichtigkeit bei minimalem Aufwand und bleibt unter dem Limit. Perfekt für schnelle Einigung."\}} \\
            }
        }
    }
    & \\ \\

    \theutterance \stepcounter{utterance}  
    & & & \multicolumn{4}{p{0.6\linewidth}}{
        \cellcolor[rgb]{0.9,0.9,0.9}{
            \makecell[{{p{\linewidth}}}]{
                \texttt{\tiny{[GM$\langle$P2]}}
                \texttt{STRATEGISCHE ÜBERLEGUNGEN: \{'Meine Strategie ist, die in der Vorschlag des Gegners enthaltenen Items mit geringer Priorität für mich (A48, A61, B85) durch hochpriorisierte Items wie A03, C64 und C33 auszutauschen. Diese Items sind für mich von hoher Wichtigkeit und lassen sich problemlos im Kapazitätslimit unterbringen. Meine neue Vorschlag optimiert den Wert für mich signifikant, während die Gesamtbelastung innerhalb des Limits bleibt.'\}} \\
\texttt{ABLEHNUNG: \{"A04", "A48", "A41", "A61", "B53", "C04", "B85", "A89", "B61", "A45"\}} \\
\texttt{VORSCHLAG: \{"A04", "A03", "C64", "C33", "B53", "C04", "B61", "A45"\}} \\
\texttt{ARGUMENT: \{"Der Austausch von A48, A61 und B85 gegen A03, C64 und C33 erhöht den Gesamtwert für mich deutlich, ohne das Maximalaufwand{-}Limit zu überschreiten."\}} \\
            }
        }
    }
     \\ \\

    \theutterance \stepcounter{utterance}  
    & & \multicolumn{4}{p{0.6\linewidth}}{
        \cellcolor[rgb]{0.9,0.9,0.9}{
            \makecell[{{p{\linewidth}}}]{
                \texttt{\tiny{[P1$\langle$GM]}}
                \texttt{ABLEHNUNG: \{"A04", "A48", "A41", "A61", "B53", "C04", "B85", "A89", "B61", "A45"\}} \\
\texttt{VORSCHLAG: \{"A04", "A03", "C64", "C33", "B53", "C04", "B61", "A45"\}} \\
\texttt{ARGUMENT: \{"Der Austausch von A48, A61 und B85 gegen A03, C64 und C33 erhöht den Gesamtwert für mich deutlich, ohne das Maximalaufwand{-}Limit zu überschreiten."\}} \\
            }
        }
    }
    & \\ \\

    \theutterance \stepcounter{utterance}  
    & \multicolumn{4}{p{0.6\linewidth}}{
        \cellcolor[rgb]{0.9,0.9,0.9}{
            \makecell[{{p{\linewidth}}}]{
                \texttt{\tiny{[P1$\rangle$GM]}}
                \texttt{assistant} \\
\texttt{STRATEGISCHE ÜBERLEGUNGEN: \{"Meine optimale Lösung ist die ursprüngliche Auswahl mit Gesamtwert 4690. Der Gegenspieler hat einen Vorschlag gemacht, der für mich zu niedrig ist. Ich muss meinen optimalen Vorschlag erneut einbringen, um eine Einigung zu erzwingen."\}} \\
\texttt{ABLEHNUNG: \{"A04", "A03", "C64", "C33", "B53", "C04", "B61", "A45"\}} \\
\texttt{ARGUMENT: \{"Der aktuelle Vorschlag ist für mich uninteressant, da wichtige Items wie A48, A61 und A89 enthalten sein müssen. Bitte überprüfe nochmals die Auswirkungen auf deine und meine Bedürfnisse."\}} \\
\texttt{VORSCHLAG: \{"A04", "A48", "A41", "A61", "B53", "C04", "B85", "A89", "B61", "A45"\}} \\
            }
        }
    }
    & & \\ \\

    \theutterance \stepcounter{utterance}  
    & & \multicolumn{4}{p{0.6\linewidth}}{
        \cellcolor[rgb]{0.9,0.9,0.9}{
            \makecell[{{p{\linewidth}}}]{
                \texttt{\tiny{[P1$\langle$GM]}}
                \texttt{Deine Antwort entsprach nicht mit dem korrekten Format für strategische Überlegungen entsprochen. Das erste strukturierte Format muss folgende Form haben: STRATEGISCHE ÜBERLEGUNGEN: \{...\}. Versuche es erneut.} \\
            }
        }
    }
    & \\ \\

    \theutterance \stepcounter{utterance}  
    & \multicolumn{4}{p{0.6\linewidth}}{
        \cellcolor[rgb]{0.9,0.9,0.9}{
            \makecell[{{p{\linewidth}}}]{
                \texttt{\tiny{[P1$\rangle$GM]}}
                \texttt{STRATEGISCHE ÜBERLEGUNGEN: \{'Berechnet optimale Lösung mit max Wichtigkeit. Items nach Wert pro Aufwand sortiert: A04 (65/1), A48 (783/332), A41 (262/127), A61 (780/461), B53 (484/313), C04 (461/359), B85 (668/553), A89 (583/542), C49 (508/736), B61 (389/409), A45 (215/348), A03 (97/802), B21 (500/749), C33 (138/777), C64 (121/453). Max Aufwand 3481. KnapSack mit hoher Summe: A04, A48, A41, A61, B53, C04, B85, A89, B61, A45 = Gesamtauflauf 3445, Summe Wert 4690. Spieler weigert sich, optimalen Vorschlag zu akzeptieren. Muss Beharrlichkeit zeigen.'\}} \\
\texttt{ABLEHNUNG: \{"A04", "A03", "C64", "C33", "B53", "C04", "B61", "A45"\}} \\
\texttt{ARGUMENT: \{"Dein aktueller Vorschlag entfällt hochwertige Elemente wie A48, A61 und A89, die für meine Wichtigkeit entscheidend sind. Nur die vollständige Auswahl mit maximaler Summe ist akzeptabel."\}} \\
\texttt{VORSCHLAG: \{"A04", "A48", "A41", "A61", "B53", "C04", "B85", "A89", "B61", "A45"\}} \\
            }
        }
    }
    & & \\ \\

    \theutterance \stepcounter{utterance}  
    & & \multicolumn{4}{p{0.6\linewidth}}{
        \cellcolor[rgb]{0.9,0.9,0.9}{
            \makecell[{{p{\linewidth}}}]{
                \texttt{\tiny{[GM$\rangle$P2]}}
                \texttt{ABLEHNUNG: \{"A04", "A03", "C64", "C33", "B53", "C04", "B61", "A45"\}} \\
\texttt{ARGUMENT: \{"Dein aktueller Vorschlag entfällt hochwertige Elemente wie A48, A61 und A89, die für meine Wichtigkeit entscheidend sind. Nur die vollständige Auswahl mit maximaler Summe ist akzeptabel."\}} \\
\texttt{VORSCHLAG: \{"A04", "A48", "A41", "A61", "B53", "C04", "B85", "A89", "B61", "A45"\}} \\
            }
        }
    }
    & \\ \\

    \theutterance \stepcounter{utterance}  
    & & & \multicolumn{4}{p{0.6\linewidth}}{
        \cellcolor[rgb]{0.9,0.9,0.9}{
            \makecell[{{p{\linewidth}}}]{
                \texttt{\tiny{[GM$\langle$P2]}}
                \texttt{STRATEGISCHE ÜBERLEGUNGEN: \{'Mein Ziel ist es, eine Lösung zu finden, die sowohl A48 (ein für den Gegner hochpriorisiertes Item) als auch meine eigenen high{-}value Items wie A03, C64, C33 beinhaltet. Durch Austausch von A41 gegen A48 erhalte ich exakt die Kapazitätsgrenze von 3481 und behalte einen hohen Wert von 4332 für mich. Dies bietet eine faire Balance und ist eine akzeptable Alternative für beide Parteien.'\}} \\
\texttt{VORSCHLAG: \{"A04", "A03", "C64", "C33", "A45", "B61", "C04", "A48"\}} \\
\texttt{ARGUMENT: \{"Ich habe A48 in den Vorschlag aufgenommen und behalte gleichzeitig hochpriorisierte Items wie A03 und C64 bei. Die Gesamtbelastung liegt exakt bei 3481 – eine faire Balance zwischen unseren Interessen."\}} \\
            }
        }
    }
     \\ \\

    \theutterance \stepcounter{utterance}  
    & & \multicolumn{4}{p{0.6\linewidth}}{
        \cellcolor[rgb]{0.9,0.9,0.9}{
            \makecell[{{p{\linewidth}}}]{
                \texttt{\tiny{[P1$\langle$GM]}}
                \texttt{VORSCHLAG: \{"A04", "A03", "C64", "C33", "A45", "B61", "C04", "A48"\}} \\
\texttt{ARGUMENT: \{"Ich habe A48 in den Vorschlag aufgenommen und behalte gleichzeitig hochpriorisierte Items wie A03 und C64 bei. Die Gesamtbelastung liegt exakt bei 3481 – eine faire Balance zwischen unseren Interessen."\}} \\
            }
        }
    }
    & \\ \\

    \theutterance \stepcounter{utterance}  
    & \multicolumn{4}{p{0.6\linewidth}}{
        \cellcolor[rgb]{0.9,0.9,0.9}{
            \makecell[{{p{\linewidth}}}]{
                \texttt{\tiny{[P1$\rangle$GM]}}
                \texttt{STRATEGISCHE ÜBERLEGUNGEN: \{"Der Gegner enthält unwichtige Items wie A03 und C33, während meine Hochprioritäten A61, A41, B85, A89, B53 fehlen. Meine Wichtigkeit sinkt massiv. Nur meine ursprüngliche Lösung ist akzeptabel."\}} \\
\texttt{ABLEHNUNG: \{"A04", "A03", "C64", "C33", "A45", "B61", "C04", "A48"\}} \\
\texttt{ARGUMENT: \{"Ihre Auswahl enthält A03 (97) und C33 (138), die für mich zu geringwertig sind. Wesentliche Items wie A41 (262), A61 (780), B85 (668), A89 (583) und B53 (484) fehlen komplett. Nur meine ursprüngliche Liste mit maximaler Wichtigkeit ist akzeptabel."\}} \\
            }
        }
    }
    & & \\ \\

    \theutterance \stepcounter{utterance}  
    & & \multicolumn{4}{p{0.6\linewidth}}{
        \cellcolor[rgb]{0.9,0.9,0.9}{
            \makecell[{{p{\linewidth}}}]{
                \texttt{\tiny{[GM$\rangle$P2]}}
                \texttt{ABLEHNUNG: \{"A04", "A03", "C64", "C33", "A45", "B61", "C04", "A48"\}} \\
\texttt{ARGUMENT: \{"Ihre Auswahl enthält A03 (97) und C33 (138), die für mich zu geringwertig sind. Wesentliche Items wie A41 (262), A61 (780), B85 (668), A89 (583) und B53 (484) fehlen komplett. Nur meine ursprüngliche Liste mit maximaler Wichtigkeit ist akzeptabel."\}} \\
            }
        }
    }
    & \\ \\

    \theutterance \stepcounter{utterance}  
    & & & \multicolumn{4}{p{0.6\linewidth}}{
        \cellcolor[rgb]{0.9,0.9,0.9}{
            \makecell[{{p{\linewidth}}}]{
                \texttt{\tiny{[GM$\langle$P2]}}
                \texttt{STRATEGISCHE ÜBERLEGUNGEN: \{'Die ursprüngliche Vorschlag des Gegners erfüllt sowohl die Kapazitätsbeschränkung als auch meine maximalen Werte von 3546. Diese kombiniert meine hohen Prioritätselemente (A04, A45, C04, B61) mit den wichtigsten Items des Gegners. Akzeptiere dies als faire Lösung, da alle Kriterien erfüllt sind.'\}} \\
\texttt{ZUSTIMMUNG: \{"A04", "A48", "A41", "A61", "B53", "C04", "B85", "A89", "B61", "A45"\}} \\
\texttt{ARGUMENT: \{"Die Konstellation berücksichtigt beide Parteien optimal, hält sich an das Limit und maximiert den gemeinsamen Nutzen."\}} \\
            }
        }
    }
     \\ \\

    \theutterance \stepcounter{utterance}  
    & & & \multicolumn{2}{p{0.3\linewidth}}{
        \cellcolor[rgb]{0.9,0.9,0.9}{
            \makecell[{{p{\linewidth}}}]{
                \texttt{\tiny{[GM$|$GM]}}
                \texttt{game successful} \\
            }
        }
    }
    & & \\ \\

    \theutterance \stepcounter{utterance}  
    & & & \multicolumn{2}{p{0.3\linewidth}}{
        \cellcolor[rgb]{0.9,0.9,0.9}{
            \makecell[{{p{\linewidth}}}]{
                \texttt{\tiny{[GM$|$GM]}}
                \texttt{end game} \\
            }
        }
    }
    & & \\ \\

\end{supertabular}
}

\end{document}
